\chapter*{要旨}
\addcontentsline{toc}{chapter}{要旨}

本研究は，災害時など安全な水資源の利用が制限される状況を想定し，
手指を介した水系感染症の感染リスクを定量的に評価することを目的とした。
特に，掌に付着する水分量およびウイルス量を実験的に測定し，
洗浄条件の違いがウイルス除去効率および感染リスク低減に及ぼす影響を検討した。

まず，被験者9名を対象として掌面積および掌体積の評価を行った。
スキャナー画像を用いた画像解析により算出した掌面積は，
おおよそ150～190\,cm$^{2}$の範囲に分布し，
左右差について統計解析を行った結果，有意差は認められなかった。
一方，掌体積については，利き手側が非利き手側より大きい傾向を示し，
統計的にも有意差が確認された（$p<0.05$）。

次に，蒸留水への浸漬後の秤量測定により，
掌表面に付着する水分量を定量した。
掌全体に付着する水分量は概ね2～4\,mLであり，
掌面積で正規化した水分付着量は
約0.01～0.02\,mL/cm$^{2}$の範囲であった。
水分付着量についても左右差は認められず，
面積補正を行うことで被験者間比較が可能であることが示された。

さらに，指標ウイルスとして大腸菌ファージMS2を用い，
手指洗浄によるウイルス除去効率を評価した。
清浄水500\,mLによる単回洗浄後のウイルス残存率は
おおよそ10～30\%であり，
洗浄後も一定量のウイルスが掌表面に残存することが確認された。
また，同一水量を複数回に分割して洗浄した場合，
初回洗浄により残存率は大きく低下したものの，
2回目以降の追加的な除去効果は限定的であった。

これらの実験結果を基に，
河川水中のノロウイルス濃度，
手指への水分付着量，
洗浄後ウイルス残存率，
および手から口への誤飲摂取割合を確率分布として設定し，
定量的微生物リスク評価（QMRA）を実施した。
モンテカルロシミュレーションの結果，
洗浄を1回行った場合の感染確率は
おおむね$10^{-4}$～$10^{-6}$のオーダーに低減することが示されたが，
感染リスクを完全にゼロにすることは困難であることが明らかとなった。

以上より，本研究は，
限られた水資源下における手指洗浄の有効性と限界を
定量的に示すとともに，
災害時における実践的な手指衛生対策の検討に
有用な基礎的知見を提供するものである。

\bigskip
\noindent
\textbf{キーワード：}
手指衛生，ウイルス付着量，大腸菌ファージMS2，
洗浄効率，定量的微生物リスク評価（QMRA）

\newpage
